\section{Основные этапы развития нелинейной оптики}\label{sec:chFiber3/fiberHist}
\fixme{напиши про немецкую женщину!}
Впервые идеи нелинейной оптики были высказаны советским учёным С. И. Вавиловым в работе <<Микроструктура света>> \cite{vavilov1950}. Обобщая свои наблюдения, относящиеся к 1920 гг., и последующие опыты, он утверждал, что нелинейность по величине световой мощности должна наблюдаться в отношении не только абсорбции, но и других оптических свойств среды: двулучепреломления, дихроизма, вращательной способности и т.д. Несмотря на то, что Вавилов указал лишь на часть физических явлений, приводящих к нелинейности среды по световому полю (изменение числа атомов или молекул, способных поглощать свет), и, как следствие, ряд  нелинейных явлений, основанных на других физических механизмах, были предсказаны и продемонстрированы позднее, \fixme{С. И. Вавилов по праву считается основоположником нелинейной оптики.}

Впервые генерация второй гармоники (ГВГ) \cite{franken1961generation} наблюдалась спустя всего год после создания первого лазера в 1960 году \cite{maiman1960stimulated}. Несмотря на то, что для эффективного протекания процесса необходимо соблюдение условий т.н. фазового синхронизма, в данной работе эффект наблюдался в его отсутствие. Затем в 1962 году А. Армстронгом и Н. Бломбергеном  были предложены три способа достижения фазового синхронизма дискретного типа или \textsf{квазисинхронизма} \cite{armstrong1962interactions}, при котором периодически контролируемо изменялась разность фаз взаимодействующих волн, а американский учёный Джордмэйн предложил использовать для достижения этой цели двулучепреломляющие кристаллы \cite{giordmaine1962mixing}. Такой тип синхронизма принято называть \textsf{традиционным}. 

Последовавший за публикацией \cite{giordmaine1962mixing} бурный рост числа исследований по генерации второй гармоники на основе традиционного синхронизма привёл к получению уже в 1963 году преобразования с эффективностью 20\%, а идеи Н. Бломбергена с сотрудниками в силу технологических трудностей создания структур для реализации квазисинхронизма были отодвинуты на второй план. Лишь в 90-х годах XX века произошёл возврат к идеям квазисинхронизма на основе кристаллов с регулярной доменной структурой. В последние годы значительная часть нелинейно-оптических преобразователей света в своей работе использует именно данный тип фазового синхронизма.

Наряду с нелинейно-оптическими процессами второго порядка, в объёмных средах и волоконных светододах изучались также и нелинейно-оптические процессы третьего порядка. \textbf{Комбинационное, или Рамановское рассеяние} было открыто независимо советским учёным Мандельштамом в кристаллах и индийским учёным Раманом~--- в жидкостях, в 1928 году. Присуждение за это открытие Нобелевской премии \fixme{в каком году} только Раману до сих пор является \fixme{предметом взаимных упрёков.} В 1962 году с использованием \fixme{рубинового} лазера был открыт эффект вынужденного комбинационного рассеяния (ВКР).
\nomenclature{ВКР}{Вынужденное комбинационное рассеяние}

Рассеяние света на акустических фононах в твёрдых телах и жидкостях было предсказано в 1914 году и экспериментально продемонстрировано в 1930 году. Вынужденное же рассеяние лазерного излучения на таких фононах было обнаружено в 1964 году и носит название \textbf{вынужденного рассеяние Мандельштама-Бриллюэна (ВРМБ)}
\nomenclature{ВРМБ}{Вынужденное рассеяние Мандельштама-Бриллюэна}

Оба вышеназванных эффекта активно исследовались в волоконных световодах в 1972 г. ССЫЛКИ, поскольку к этому моменту уровень оптических потерь в световодах приблизился к фундементальному ограничению, обусловленному процессом Рэлеевского рассеяния, и нелинейно-оптические эффекты третьего порядка стали проявляться особенно ярко. ХЫ.

Нелинейно-оптические эффекты третьего порядка, обусловленные зависимостью показателя преломления среды от величины светового поля и не зависящение от фазовых соотношений взаимодействующих волн~--- фазовая кросс-модуляция (ФКМ) и фазовая самомодуляция (ФСМ)~--- были открыты в \fixme{каком-то году}
\nomenclature{ФСМ}{Фазовая кросс-модуляция}
\nomenclature{ФКМ}{Фазовая самомодуляция}

 
\fixme{Что-то написать про историю нелинейщины третьего порядка}